% ---------------------------------------------------------------------------
% 10 Conclusion and closing statements
% ---------------------------------------------------------------------------

\section{Conclusion}

This study gives a split answer to the question of encrypted DNAGPT inference. Client-assisted CKKS
executes all twelve released blocks and the task head at the $103$-token prompt from encrypted
embedded vectors. For the one input evaluated, the decrypted label matches the frozen reference and
its margin has $8.56\times10^{-9}$ relative error against the $4\times10^{-2}$ predeclared tolerance.
Graph feasibility is established for the complete released computation; numerical generality
beyond this input is not. Client refresh bounds depth by the longest encrypted segment,
device-wide GPU memory used stays at $9839$~MiB across the blocks, and no homomorphic bootstrap is
required. A stage-bounded plaintext cache separately prevents host-memory accumulation.

Cost is the constraint that remains. One encrypted classification of one $600$-base-pair sequence
was measured at $1.86$~hours on an A100 GPU node, which forecloses interactive use in the evaluated
form. It is a single-process, single-input, single-execution result rather than a service-rate or
network measurement, and sampled GPU utilization peaking at $51\%$ indicates unclaimed headroom
without identifying its cause. The two verdicts are therefore separate, and both are informative:
complete-model feasibility is established for the executed case, practicality is not, and the
obstacle between them is engineering throughput on the provider's encrypted path rather than
correctness, multiplicative depth, or accelerator memory.

\subsection*{Future work}

Further empirical work would require additional prompts, repeated timings, and a paired
pre-optimization measurement; none is inferred from the present run. Systems work should extend the
encode-once strategy from weight diagonals to mask plaintexts. Distinct transforms across the
packed activation copies are the larger circuit opportunity. Before mask and copy-repair overhead,
the derived schedule for that layout projects a reduction of roughly $60\%$ in the block's total
ciphertext--plaintext count. A thread-safe cryptographic context would permit concurrent independent
client boundaries. Protocol closure also requires a noise-flooding budget for re-encryption. A
deployment study must add serialized network transport and private token-index lookup before
measuring the complete service.

\section*{Ethics statement}

This study involved no human participants, specimens, or private clinical records. It used public
reference datasets with recorded provenance. The work makes no clinical-performance claim and
does not present the prototype as a deployable privacy guarantee beyond its stated threat model.

\section*{Data availability statement}

All evaluation datasets are public. Their original sources, versions, preprocessing, and any
recovery route used are recorded in the repository's data-provenance document. Sequence data,
model weights, and cryptographic keys are not redistributed.

\section*{Code availability statement}

The evaluation harnesses, encrypted-evaluation implementation, environment definitions, and the
evidence record for the paper are maintained at
\url{https://github.com/Georgakopoulos-Soares-lab/private_genomic_dnagpt}. No repository-level
reuse license has yet been assigned.
